\documentclass[a4paper,11pt]{article}
\usepackage[utf8]{inputenc}
\usepackage[T1]{fontenc}
\usepackage[french]{babel}
\usepackage[left=2cm,right=2cm,top=2cm,bottom=2cm]{geometry}
\usepackage{xcolor}
\usepackage{hyperref}
\usepackage{fontawesome5}
\usepackage{parskip}

% --- Couleurs RATP ---
\definecolor{primary}{RGB}{0, 100, 60}
\hypersetup{colorlinks=true, urlcolor=primary}

\begin{document}

% --- EXPÉDITEUR ---
\noindent
\begin{minipage}[t]{0.5\textwidth}
    \textbf{Loïc Aron MBASSI EWOLO}\\
    Élève-Ingénieur Bac+5 -- Double diplôme ENSEA\\[3pt]
    \faMapMarker\ Cergy, France\\
    \faPhone\ (+33) 7 59 52 33 98\\
    \faEnvelope\ \href{mailto:loicaron.mbassiewolo@ensea.fr}{loicaron.mbassiewolo@ensea.fr}
\end{minipage}
\hfill
\begin{minipage}[t]{0.4\textwidth}
    \raggedleft
    Cergy, le 20 mars 2026
\end{minipage}

\vspace{1.2cm}

\hfill
\begin{minipage}[t]{0.5\textwidth}
    \raggedleft
    \textbf{Groupe RATP}\\
    Direction des Systèmes d'Information\\
    Data Factory\\
    Paris, France
\end{minipage}

\vspace{1cm}

\noindent\textbf{Objet :} Candidature pour l'alternance ML Engineer IA Générative (DSI/FAB)

\vspace{0.8cm}

Madame, Monsieur,

Cinq agents IA collaboratifs orchestrés via LangChain et Google ADK, un pipeline RAG complet (vectorisation, recherche sémantique, augmentation contextuelle), le tout déployé en production avec FastAPI et Docker : c'est le système multi-agents que j'ai conçu dans le cadre d'un projet de recherche en 2024-2025. En découvrant l'offre d'alternance ML Engineer au sein de votre Data Factory, j'ai immédiatement reconnu les enjeux techniques que vous décrivez, et c'est ce qui m'amène à vous adresser ma candidature.

L'expérience la plus directement transposable à vos missions est ce système multi-agents agricole : orchestration de LLMs (Gemini) pour des recommandations contextualisées, pipeline RAG avec gestion du chunking, du prompting et de l'évaluation de la pertinence des réponses, intégration d'APIs REST externes et résolution de conflits entre agents. En parallèle, j'ai participé au développement de chatbots intelligents RAG sur le projet Snappy (messagerie d'entreprise), avec recherche sémantique par embeddings vectoriels et enrichissement contextuel des réponses à partir de documents d'organisation. Ces deux projets m'ont appris les mécanismes concrets de mise en production de solutions RAG, de l'ingestion documentaire à l'évaluation des réponses.

Mon stage de recherche à l'INRIA Saclay (Équipe TRIBE, Institut Polytechnique de Paris, Avril-Août 2026) consolide mon profil sur le plan du NLP à grande échelle : analyse de textes par BERT/Transformers, conception de pipelines Python robustes, et modélisation de scores interprétables. Avant cela, mon passage au MILA (Institut Québécois d'IA, été 2025) m'a confronté à l'étude du comportement des réseaux de neurones profonds, un socle essentiel pour comprendre les forces et limites des LLMs que je serais amené à exploiter chez vous.

Je suis disponible dès septembre 2026 pour intégrer la Data Factory. Contribuer à l'industrialisation de vos assistants IA dans un contexte aussi concret que le transport francilien, au service de milliers d'agents et de voyageurs, représente pour moi une opportunité de donner une dimension opérationnelle directe à mes travaux sur les LLMs et les systèmes multi-agents. Je serais ravi d'en discuter lors d'un entretien.

Je vous prie d'agréer, Madame, Monsieur, l'expression de mes salutations distinguées.

\vspace{0.8cm}

\textbf{Loïc Aron MBASSI EWOLO}\\
\textit{Élève-Ingénieur ENSEA / Polytechnique Yaoundé}

\end{document}
